\chapter{The Panel Fights Back}
\label{ch:panel}

\section{What was built}

The first piece of competition furniture: \texttt{/World/IDMOPanel} --- a kinematic board
carrying a spring-return push button (prismatic joint, 8\,mm travel, drive-as-spring) and a
toggle switch (revolute, $\pm25^\circ$) --- generated \emph{procedurally} by
\texttt{ros2\_ws/isaac/build\_panel.py}. Procedural matters: a scripted panel is a
randomizable panel, and randomization is the P1 benchmark's core requirement. The build took
one evening; making the \emph{arm and panel coexist} took the rest of it, and taught more
than the build did.

\section{Interactive elements are tiny robots}

A button is a 1-DoF machine: body + prismatic joint + drive whose stiffness plays the return
spring; ``pressed'' is a joint-position threshold, readable from simulation state exactly like
an arm joint. The switch is the same idea rotated: revolute + hard limits + light damping so
it rests where flipped. Success detection for the whole IDMO skill family will be \emph{joint
and pose queries on the panel}, not vision tricks --- the simulator gives us ground truth for
free.

\section{The debugging ledger}

\begin{lesson}
\textbf{A wall is not a rigid body on a joint.} Panel v0 anchored the board with a
\texttt{FixedJoint} carrying \emph{no anchor frames}; on Play, PhysX welded board-origin to
world-origin with whatever relative pose it derived --- the board pivoted into a fallen-sail
pose. v0.1 rule: immovable scenery is a \emph{kinematic rigid body} (infinite mass, still
collidable), no joint to the world at all. Joints are for things that move.
\end{lesson}

\begin{lesson}
\textbf{Joint anchors live in scaled local frames.} \texttt{UsdPhysics} joint
\texttt{localPos} is expressed in the body's local frame \emph{and multiplied by the prim's
scale}. Anchors tuned under the earlier half-scale bug pointed nowhere after the scale fix ---
elements snapped to phantom seats. v0.1 writes anchors in pre-scale units with the arithmetic
in comments, so the next reader can check the numbers instead of trusting them.
\end{lesson}

\begin{lesson}
\textbf{The controller never forgets, and now the workspace has furniture.} The JTC holds its
last commanded pose indefinitely. The moment a solid panel occupied the arm's old sweep path,
every Play cycle made the arm lunge \emph{through} the board (PhysX resolving interpenetration
looks exactly like ``an impossible state''). Disciplines adopted: park the arm at a home pose
before Stop/Play cycles; audit stored trajectories when the scene changes; and ---
the real fix, next chapter --- collision-aware planning, which is what MoveIt is \emph{for}.
\end{lesson}

\begin{lesson}
\textbf{Strict tolerances block rescues too.} Escaping the panel kept aborting: the arm
physically could not track the commanded path while grinding against geometry, and our
freshly-installed 0.2\,rad path tolerance --- working exactly as designed --- kept pulling the
plug mid-rescue. Escape choreography that worked: rotate \emph{away} from the obstacle sector
first (free axis), \emph{then} fold home. A safety system that also constrains recovery is a
feature to design around, not a bug to loosen.
\end{lesson}

\begin{lesson}
\textbf{Kill your processes by evidence, not by faith.} The ``dead'' control node wasn't; a
relaunch stacked a second controller manager onto the first (symptom: \emph{``cannot be
configured from `active' state''}), and three zombie \texttt{robot\_state\_publisher}s from
successive launches were still publishing. Rules: after any kill, verify with
\texttt{pgrep -fa}; before any launch, know what is already running. Dueling controller
managers produce Kafkaesque symptoms --- goals accepted by one, hardware owned by another.
\end{lesson}

\begin{lesson}
\textbf{Springs need dampers: read the asset's numbers.} The arm drifted between goals,
overshot every motion, and failed tolerances on a different joint each attempt. The
$\times10$-and-print diagnostic revealed the UR5e asset ships position drives with stiffness
$\sim10^4$ but damping $\sim0.06$--$4$ --- a nearly undamped oscillator. Setting damping to
$5\%$ of stiffness (thousands, not units) gave crisp, dead-beat motion and the first
\texttt{SUCCEEDED} home. Meta-rule: when tracking is mysteriously sloppy, print the drive
gains before theorizing --- and prefer diagnostics that \emph{report} state over ones that
just mutate it.
\end{lesson}

\section{Where this leaves us}

Arm parked at a tuned home; panel functional (button cycles under its drive, switch rests at
its detents); every joint drive in the scene now has explicit, recorded gains. The arm's tool
flange spent part of the evening accidentally resting on the red button --- an unplanned
preview of the next chapter, where MoveIt plans a deliberate, collision-aware press of the
same button. First accidental contact: 5~July. First intentional press: next session's goal.
